% !TEX TS-program = xelatex
% !TEX encoding = UTF-8 Unicode
% !Mode:: "TeX:UTF-8"

\documentclass{resume}
\usepackage{zh_CN-Adobefonts_external} % Simplified Chinese Support using external fonts (./fonts/zh_CN-Adobe/)
%\usepackage{zh_CN-Adobefonts_internal} % Simplified Chinese Support using system fonts
\usepackage{linespacing_fix} % disable extra space before next section
\usepackage{cite}


% \documentclass{resume}
% \usepackage{graphicx}
% \usepackage{zh_CN-Adobefonts_external}
% \usepackage{linespacing_fix}
% \usepackage{cite}



\begin{document}
\pagenumbering{gobble} % suppress displaying page number

% \name{杜靖怡}

% \contactInfo{(+86) 185-3792-2003}{dujingyi7714@163.com}
% {GitHub @hijiangtao}


\noindent
\begin{tabular*}{\textwidth}{@{\extracolsep{\fill}} p{0.72\textwidth} p{0.22\textwidth}}
    \begin{minipage}[c][3.35cm][c]{0.72\textwidth}
        \centering
        {\Huge \textbf{杜靖怡}}
        
        \vspace{0.36cm}
        \normalsize 电话：185-3792-2003 \quad 邮箱：dujingyi7714@163.com
    \end{minipage}
    &
    \begin{minipage}[t][3.35cm][t]{0.22\textwidth}
        \centering
        \vspace{0.75cm}
        % image removed for test
    \end{minipage}
\end{tabular*}

\vspace{0.12cm}

% % \section{\faGruationCap\ 教育背景}
% \section{教育背景}
% \datedsubsection{\textbf{中国科学院大学},计算机应用技术,\textit{在读硕士研究生}}{2015.9 - 2018.6}
% \ \textbf{排名11/133(前10\%)},中国科学院大学学业奖学金(2次),IEEE Student member,预计2018年6月毕业
% \datedsubsection{\textbf{北京理工大学},软件工程,\textit{工学学士}}{2011.9 - 2015.6}
% \ \textbf{排名2/62(前5\%)},国家励志奖学金,人民奖学金(7次),科技竞赛奖(2次),北京市普通高等学校优秀毕业生,北京理工大学优秀毕业生,软件学院金牌毕业生,优秀团员/优秀学生(5次)
% % \end{itemize}




\section{教育经历}

{
\renewcommand{\arraystretch}{1.25}
\begin{tabular*}{\textwidth}{@{\extracolsep{\fill}} l l l r}
\textbf{吉林大学(985)} & 统计学 & 硕士在读 & 2025.09 - 2028.06 \\
\textbf{湘潭大学(双一流)} & 统计学 & 理学学士 & 2021.09 - 2025.06 \\
\end{tabular*}
}

\vspace{-0.6em}

\begin{itemize}
    \item \textbf{GPA}: 3.34/4.0 \hfill \textbf{CET6}: 520 \hfill \textbf{雅思}: 6.5
    \item \textbf{奖项/荣誉}: 全国大学生数学建模竞赛省一、统计建模大赛省二、校级三好学生、湖南省优秀毕业生
\end{itemize}

\section{项目经历}
\datedsubsection{\textbf{基于多行为序列建模的 OTTO 电商推荐与用户行为预测}}{2026.06-2026.07}
\begin{itemize}

  \item 基于 OTTO 电商行为数据构建 clicks/carts/orders 多目标推荐任务，处理约 6014 万条训练日志、208 万条验证日志、45 万个 session，完成数据清洗、行为编码与验证集构造。
  \item 设计 Item-to-Item 共现召回策略，融合 clicks、carts-orders、buy2buy 三类共现矩阵，结合 行为权重 1:5:4、时间窗口和时间衰减 生成候选商品。
  \item 构造商品侧、session侧、session-aid 交互侧共 37 维排序特征，生成约 6755 万条候选样本，并通过分片存储和负采样缓解内存压力与样本不平衡问题。
  \item 使用 XGBoost rank:pairwise + GroupKFold 训练三类排序模型，共训练 15 个模型，本地验证集达到 Weighted Recall@20 = 0.5410，其中 orders 任务 Recall@20 = 0.6234
\end{itemize}

\datedsubsection{\textbf{量子卷积特征提取的混合经典-量子图像分类研究}}{2025.11-2025.12}
\begin{itemize}
  \item 围绕量子卷积特征提取在图像分类任务中的应用开展研究，基于 Fashion-MNIST 构建 Classical CNN 与 Quantum-enhanced CNN 对照实验，完成图像预处理、量子特征生成、模型训练与结果分析。
  \item 基于 PyTorch 自定义 CNN 模型，搭建两层 Conv2d-BatchNorm-ReLU-MaxPool 结构，使用 CrossEntropyLoss 与 Adam 优化器完成训练，记录 train/test loss 和 accuracy 曲线。
  \item 分析量子卷积层在图像特征提取中的作用及计算开销，并讨论随机量子电路、编码方式和经典-量子对照实验中的控制变量问题。
\end{itemize}

\datedsubsection{\textbf{北京格灵深瞳信息技术有限公司 | DeepGlint},Web开发工程师}{2015.7-2015.9}
\begin{itemize}
  \item \textbf{独立负责MUSE部门的可视化组件研发。}与平台研发、设计协作完成 DeepGlint Developer 平台可视化图表组件的集成开发，符合完全定制化渲染、响应式布局与实时更新等特点
  \item 利用 D3+Vue+WebGL(Three.js) 尝试实现三维空间的人群移动可视化
\end{itemize}




% \section{\faCogs\ IT 技能}
\section{专业技能}
% increase linespacing [parsep=0.5ex]
\begin{itemize}[parsep=0.2ex]
  \item \textbf{编程与数据处理}: 熟悉 Python，具备数据清洗、特征工程、实验分析和模型评估能力；了解 SQL、Git、Linux 基础操作；能够熟练使用Codex、Claude Code等AI开发工具。
  \item \textbf{机器学习与深度学习}: 掌握 决策树、随机森林、GBDT、XGBoost 等机器学习算法，熟悉交叉验证、特征选择、类别不平衡处理和模型调参方法；熟悉 PyTorch 框架，了解 MLP、CNN、RNN/LSTM、Transformer 等深度学习模型结构，具备复现和调试能力。

  \item \textbf{数理统计基础}:具备扎实的概率论、数理统计、回归分析、时间序列分析和运筹优化基础；熟悉假设检验、ARIMA、非线性规划等方法，具备统计建模、预测分析和业务问题抽象能力。
\end{itemize}

% \section{\faHeartO\ 项目/作品摘要}
% \section{项目/作品摘要}
% \datedline{\textit{An Integrated Version of Security Monitor Vis System}, https://hijiangtao.github.io/ss-vis-component/ }{}
% \datedline{\textit{Dark-Tech}, https://github.com/hijiangtao/dark-tech/ }{}
% \datedline{\textit{融合社交网络数据挖掘的电视节目可视化分析系统}, https://hijiangtao.github.io/variety-show-hot-spot-vis/}{}
% \datedline{\textit{LeetCodeOJ Solutions}, https://github.com/hijiangtao/LeetCodeOJ}{}
% \datedline{\textit{Info-Vis}, https://github.com/ISCAS-VIS/infovis-ucas}{}




%% Reference
%\newpage
%\bibliographystyle{IEEETran}
%\bibliography{mycite}
\end{document}
